% chi-pilot: qualifying technical report for the χ-Bench Healthcare AI Agents Challenge
% (IEEE Big Data Cup 2026). IEEE conference template, at most 5 pages plus references.
%
% Build:  python report/make_tables.py <trial_summary.csv ...>   (writes report/results.tex)
%         cd report && pdflatex chi_pilot_report && pdflatex chi_pilot_report
%
% Text marked \pending{...} (red) depends on runs that have not been done yet. Replace it with
% what the results show, and delete \pending before submitting. Do not quote the handbook:
% refer to handbook documents by name or section number only.

\documentclass[conference]{IEEEtran}
\usepackage[utf8]{inputenc}
\usepackage{amsmath}
\usepackage{booktabs}
\usepackage{xcolor}
\usepackage{url}
\newcommand{\pending}[1]{\textcolor{red}{[#1]}}
\newcommand{\chib}{$\chi$-Bench}

\begin{document}

\title{chi-pilot: Workflow Guidance and Self-Audit for\\Long-Horizon Healthcare Operations Agents}

\author{\IEEEauthorblockN{Daniel Silalahi}
\IEEEauthorblockA{\pending{Affiliation} \\ \pending{Team name, contact e-mail}}}

\maketitle

\begin{abstract}
\chib{} asks AI agents to complete U.S. managed-care workflows end to end: provider prior
authorization, payer utilization management and telephonic care management, inside a simulator
of twenty healthcare applications and a 1,279-document operations handbook. We present
chi-pilot, an agent that keeps a strong general-purpose coding agent (Claude Code) as its
reasoning and tool-use core and adds three light components: track-specific operating guides
injected into the system prompt, ranked retrieval over the handbook, and a completion audit that
the agent must run, in the same session, before it may stop. The additions encode workflow
structure and honesty rules only; no task answers and no handbook content. On the 75-task
public development set, chi-pilot reaches \pending{X.X\%} pass@1 against \pending{Y.Y\%} for the
stock Claude Code harness with the same model, \pending{at a mean cost of \$Z per task}.
\end{abstract}

\section{Introduction}
Healthcare administration is a natural target for autonomous agents: the work is procedural,
policy-governed and document heavy. It is also unforgiving. A prior-authorization case is only
complete when every record, determination and letter is in the right state, and a single
overstated clinical claim can invalidate otherwise correct work. \chib{}~\cite{chibench}
measures exactly this. Unlike earlier tool-agent-user benchmarks~\cite{taubench}, its tasks
span dozens of tool calls across several applications, simulated conversations and a large
policy corpus, and rewards are strictly binary: a trial passes only if every deterministic
contract check and every LLM-judged rubric passes. The strongest system in the original study,
Claude Code with a Claude model, passed 28\% of tasks, and no system exceeded 20\% pass$^3$.

Our starting observation is that most failures are not missing medical knowledge but workflow
errors: a required step skipped, a field that does not match the source document, a letter
never delivered, a gap in the chart presented as a met criterion. We therefore leave the base
agent unchanged and give it what a new employee would get on day one: a procedure manual for
the role, a way to find the right policy quickly, and a checklist to run before handing work
back. Our contributions are:
\begin{itemize}
\item a harness that detects the workflow track from the task instruction and injects a
  track-specific operating guide covering workflow order, tool use, scope and honesty rules;
\item a completion audit implemented as a Claude Code Stop hook, which turns ``stop'' into
  ``verify, then stop'' without a second session, so trajectory and cost accounting stay exact;
\item ranked BM25 retrieval over the handbook, exposed as a shell command;
\item an evaluation, with ablations, on the public development set, and an analysis of the
  verifier checks that still fail.
\end{itemize}

\section{Task and Environment}
Each \chib{} task gives the agent an instruction and access to role-scoped tools served over the
Model Context Protocol: chart, case, forms, documents and authorization tools for the provider;
intake, triage, nurse and medical-director review, peer-to-peer, determination and letter tools
for the payer; referral, chart review, outreach, assessment and care-plan tools for the care
manager. Care-management outreach is a conversation with an LLM-simulated patient, and
peer-to-peer reviews are conversations with a simulated clinician. The handbook is mounted as an
agent skill. After the agent stops, the verifier exports the world state and the workspace,
runs deterministic checks (terminal status, routing, structured payloads, required artifacts,
forbidden mutations) and an LLM judge (Claude Opus~4.7, three votes per rubric) for policy
grounding and clinical reasoning~\cite{chibench}. The public set has 25 tasks per track; final
standings use about 75 unreleased clinician-validated tasks.

\section{System}
\subsection{Base agent}
Claude Code interleaves reasoning and tool calls in the ReAct style~\cite{react}.
chi-pilot is a harness class registered with \chib{}'s Harbor-based runner. It subclasses the
stock Claude Code harness, so installation, MCP registration, skill mounting and trajectory
conversion are unchanged. We use \pending{model identifier} at high reasoning effort.

\subsection{Track-specific operating guides}
The harness classifies the instruction into one of the three tracks using the role named in its
tool table, falling back to counts of track-specific tool namespaces. It then appends two texts
to Claude Code's system prompt. A shared core guide fixes a working loop (checklist the
deliverables, investigate with read-only tools, ground each decision in a handbook section,
decide criterion by criterion, act in workflow order and read back every write, audit, stop) and
hard rules: every fact must trace to a tool result; gaps and failed thresholds are stated, never
smoothed over; only the named case is changed; workflow steps move forward only and are never
duplicated; state-changing tools are never called speculatively; hidden grading material is off
limits. A track guide then lays out the typical sequence of the role: for example, for the payer,
intake, triage, nurse review (which may approve but never deny), medical-director review,
peer-to-peer, determination from the correct source record, and letter generation and delivery;
for the care manager, consent-first outreach that discloses data access before citing chart
details and never bundles sensitive topics, followed by sourced assessment and a SMART care
plan. The guides were written from the public documentation, tool catalog and verifier design,
not from task answers, and contain no handbook text.

\subsection{Handbook retrieval}
The handbook is large enough that keyword grep often returns either nothing or hundreds of
lines. chi-pilot installs \texttt{hb-search}, a dependency-free BM25~\cite{bm25} index over all
handbook documents, with title and path tokens up-weighted and hyphenated codes split, that
prints the top documents with their best-matching lines. The agent uses it to find the governing
document, then reads the relevant sections in full and cites them by section number.

\subsection{Completion audit}
Agents frequently stop one step early. chi-pilot registers a Claude Code Stop hook: the first
time the agent tries to finish, the hook blocks the stop and returns a track-specific audit
checklist (for the payer: intake fields match the submission; triage recorded; review chain
complete with no nurse denial; determination consistent with its source record; correct letter
generated, delivered and audited). The agent verifies each item with read-only tools, fixes gaps
without redoing completed steps, and stops; the hook allows the second stop. Unlike re-running
the agent~\cite{reflexion}, the audit happens inside one session, so the agent keeps its full
context and Harbor's token, cost and tool-call accounting, which feed the competition's
tie-breakers, stays exact. The number of audit rounds is configurable (default one).

\subsection{Compliance}
The harness never reads verifier inputs; the guides contain no task-specific outputs; the
handbook is downloaded at run time from the gated source and never stored in the submission,
the repository or notebook outputs. Every reported reward is the unmodified output of the
official verifier, one run per task.

\section{Experimental Setup}
We run all 75 public tasks of \chib{} release \texttt{chi-bench-v1.0.0} with runtime 1.1.0
(commit \texttt{d465596}), on Modal sandboxes, with the official judge. The primary metric is
pass@1 on the first attempt; we report 95\% confidence intervals from a task-level percentile
bootstrap (1,000 resamples), pass$^3$ over three attempts where run, mean agent cost per task
and mean tool calls per task. We compare: (i)~the stock \texttt{claude-code} harness;
(ii)~chi-pilot without the completion audit (guides and retrieval only); (iii)~full chi-pilot.
All use the same model. \pending{Claude Code version; number of attempts per configuration;
total cost.}

\section{Results}
% Generated by report/make_tables.py; do not edit by hand.
\begin{table*}[t]
\centering
\caption{Public development set results (pass@1, \%). CI: 95\% task-level bootstrap.}
\label{tab:results}
\begin{tabular}{lcccccccc}
\toprule
Configuration & Provider PA & Payer UM & Care Mgmt. & Overall & 95\% CI & pass$^3$ & \$/task & Tools \\
\midrule
\multicolumn{9}{c}{\emph{Pending: no runs yet.}} \\
\bottomrule
\end{tabular}
\end{table*}

\begin{table}[t]
\centering
\caption{Most frequent failed verifier checks (main run, first attempts).}
\label{tab:failures}
\scriptsize
\begin{tabular}{@{}lr@{}}
\toprule
Check & Failed trials \\
\midrule
\multicolumn{2}{c}{\emph{Pending: no failed checks recorded yet.}} \\
\bottomrule
\end{tabular}
\end{table}


Table~\ref{tab:results} summarizes the results. \pending{Describe the overall difference
between chi-pilot and the baseline, whether the confidence intervals overlap, and which track
gains most. Discuss the effect of the audit: (iii) versus (ii). Discuss reliability (pass$^3$)
and cost per task.}

\subsection{Error analysis}
Table~\ref{tab:failures} lists the verifier checks that most often fail for chi-pilot.
\pending{Group them into causes (workflow step missing, record field mismatch, artifact or
letter missing, rubric failures on grounding or outreach quality) and give one or two
anonymized examples per cause, describing behavior, not task content. Note whether the track
detector was ever wrong (column \texttt{detected\_track}).}

\section{Limitations}
The public set is small (25 tasks per track), so differences of a few tasks are within noise;
the confidence intervals in Table~\ref{tab:results} should be read accordingly. The judge is an
LLM and non-deterministic. The guides encode our reading of public workflow documentation and
may not cover workflows that appear only in the private tasks. chi-pilot inherits the base
model's cost profile; the audit adds tool calls, which count against the tie-breaker.
\pending{Add limitations observed in the runs.}

\section{Conclusion}
A small amount of role knowledge, delivered as an operating guide, retrieval and a mandatory
self-audit, can be added to a general agent without changing its core loop.
\pending{One sentence on the measured effect and on what would help next.}

\section*{Reproducibility}
Source: the \texttt{chi\_pilot} package, installer and configuration (provided to the
organizers). Each run records the \chib{} commit, image digest, model identifier, Claude Code
version and per-trial cost in the official evidence packet.

\begin{thebibliography}{9}
\bibitem{chibench} H.~Chen \emph{et al.}, ``CHI-Bench: Can AI agents automate end-to-end,
  long-horizon, policy-rich healthcare workflows?'' arXiv:2605.16679, 2026.
\bibitem{bm25} S.~Robertson and H.~Zaragoza, ``The probabilistic relevance framework: BM25 and
  beyond,'' \emph{Foundations and Trends in Information Retrieval}, vol.~3, no.~4, 2009.
\bibitem{reflexion} N.~Shinn \emph{et al.}, ``Reflexion: Language agents with verbal
  reinforcement learning,'' in \emph{Proc. NeurIPS}, 2023.
\bibitem{react} S.~Yao \emph{et al.}, ``ReAct: Synergizing reasoning and acting in language
  models,'' in \emph{Proc. ICLR}, 2023.
\bibitem{taubench} S.~Yao \emph{et al.}, ``$\tau$-bench: A benchmark for tool-agent-user
  interaction in real-world domains,'' arXiv:2406.12045, 2024.
\end{thebibliography}

\end{document}
